\chapter{2006年上海卷（春）}
\section{填空题}
\begin{problem}
计算: \(\lim_{n\to \infty}\frac{3n - 2}{4n + 3} =\)
\end{problem}

\begin{problem}
方程 \( \log_{3}(2x-1)=1 \) 的解 x=\fillinblank{} \fillinblank{}.
\end{problem}

\begin{problem}
函数 \( f(x)=3x+5, x\in[0,1] \) 的反函数 \( f^{-1}(x)= \)  \fillinblank{}.
\end{problem}

\begin{problem}
不等式 \(\frac{1 - 2x}{x + 1} > 0\) 的解集是 \fillinblank{} \fillinblank{}.
\end{problem}

\begin{problem}
已知圆 \(C: (x + 5)^2 + y^2 = r^2 (r > 0)\) 和直线 \(l: 3x + y + 5 = 0\). 若圆 \(C\) 与直线 \(l\) 没有公共点，则 \(r\) 的取值范围是 \fillinblank{} \fillinblank{}.
\end{problem}

\begin{problem}
已知函数 \( f(x) \) 是定义在 \((- \infty, +\infty)\) 上的偶函数. 当 \( x \in (-\infty, 0) \) 时，\( f(x) = x - x^4 \)，则当 \( x \in (0, +\infty) \) 时，\( f(x) = \)  \fillinblank{}.
\end{problem}

\begin{problem}
电视台连续播放6个广告，其中含4个不同的商业广告和2个不同的公益广告，要求首尾必须播放公益广告，则共有种不同的播放方式(结果用数值表示) \fillinblank{}.
\end{problem}

\begin{problem}
正四棱锥底面边长为 4，侧棱长为 3，则其体积为 \fillinblank{} \fillinblank{}.
\end{problem}

\begin{problem}
在 \(\triangle ABC\) 中，已知 \(BC = 8, AC = 5\)，三角形面积为 12，则 \(\cos 2C =\)  \fillinblank{}.
\end{problem}

\begin{problem}
若向量 \(\bs{a}, \bs{b}\) 的夹角为 \(150^{\circ}\)，\(|\bs{a}| = \sqrt{3}\)，\(|\bs{b}| = 4\)，则 \(|2\bs{a} + \bs{b}| = \)  \fillinblank{}.
\end{problem}

\begin{problem}
已知直线 l 过点 \( P(2,1) \) ，且与 x 轴、y 轴的正半轴分别交于 A、B 两点，O 为坐标原点，则三角形 OAB 面积的最小值为 \fillinblank{} \fillinblank{}.
\end{problem}

\begin{problem}
同学们都知道，在一次考试后，如果按顺序去掉一些高分，那么班级的平均分将降低; 反之，如果按顺序去掉一些低分，那么班级的平均 分将提高. 这两个事实可以用数学语言描述为: 若有限数列 \(a_1, a_2, \cdots, a_n\) 满足 \(a_1 \leqslant a_2 \leqslant \cdots \leqslant a_n\)，则 (结论用数学式子表示) \fillinblank{}.
\end{problem}

\section{单选题}
\begin{problem}
抛物线 \( y^{2}=4x \) 的焦点坐标为（\quad）

\choices
    {\((0,1)\).}
    {\((1,0)\).}
    {\((0,2)\).}
    {\((2,0)\).}

\end{problem}

\begin{answer}
B
\end{answer}
\begin{problem}
若 \(a, b, c \in \R, a > b\)，则下列不等式成立的是（\quad）

\choices
    {\(\frac{1}{a} < \frac{1}{b}\).}
    {\(a^2 > b^2\).}
    {\(\frac{a}{c^2 + 1} >\frac{b}{c^2 + 1}\).}
    {\(a|c| > b|c|\).}

\end{problem}

\begin{answer}
C
\end{answer}
\begin{problem}
若 \(k \in \R\)，则 “\(k > 3\)” 是“方程 \(\frac{x^2}{k - 3} - \frac{y^2}{k + 3} = 1\) 表示双曲线”的（\quad）

\choices
    {充分不必要条件.}
    {必要不充分条件.}
    {充要条件.}
    {既不充分也不必要条件.}
\end{problem}

\begin{problem}
若集合 \(A = \left\{y \mid y = x^{\frac{1}{3}}, -1 \leqslant x \leqslant 1\right\}\)，\(B = \left\{y \mid y = 2 - \frac{1}{x}, 0 < x \leqslant 1\right\}\)，则 \(A \cap B\) 等于 ……（\quad）

\choices
    {\((-\infty, 1]\).}
    {\([-1, 1]\).}
    {\(\varnothing\).}
    {\(\{1\}\).}

\end{problem}

\begin{answer}
B
\end{answer}
\section{解答题}
\begin{problem}
在长方体 \( ABCD-A_{1}B_{1}C_{1}D_{1} \) 中，已知 \( DA = DC = 4, DD_{1} = 3 \) ，求异面直线 \( A_{1}B \) 与 \( B_{1}C \) 所成角的大小（结果用反三角函数值表示）.
\end{problem}

\begin{problem}
已知复数 w 满足 \( w - 4 = (3 - 2w) \) i (i 为虚数单位), \( z = \frac{5}{w} + |w - 2| \) ，求一个以 z 为根的实系数一元二次方程.
\end{problem}

\begin{problem}
已知函数 \( f(x) = 2\sin \left(x + \frac{\pi}{6}\right) - 2\cos x \)，\( x \in \left[\frac{\pi}{2}, \pi\right] \). (1) 若 \( \sin x = \frac{4}{5} \) ，求函数 \( f(x) \) 的值； (2) 求函数 \( f(x) \) 的值域.
\end{problem}

\begin{problem}
学校科技小组在计算机上模拟航天器变轨返回试验. 设计方案如图: 航天器运行(按顺时针方向)的轨迹方程为 \(\frac{x^2}{100} + \frac{y^2}{25} = 1\)，变轨(即航天器运行轨迹由椭圆变为抛物线)后返回的轨迹是以 \(y\) 轴为对称轴、\(M\left(0, \frac{64}{7}\right)\) 为顶点的抛物线的实线部分，降落点为 \(D(8,0)\). 观测点 \(A(4,0)\)、\(B(6,0)\) 同时跟踪航天器.
(1) 求航天器变轨后的运行轨迹所在的曲线方程; (2) 试问: 当航天器在 x 轴上方时，观测点 A、B 测得离航天器的距离分别为多少时，应向航天器发出变轨指令?
\end{problem}

\begin{problem}
设函数 \( f(x) = |x^2 - 4x - 5| \). (1) 在区间 \( [-2,6] \) 上画出函数 \( f(x) \) 的图象； (2) 设集合 \(A = \{x \mid f(x) \geqslant 5\}\)，\(B = (-\infty, -2] \cup [0, 4] \cup [6, +\infty)\). 试判断集合 \(A\) 和 \(B\) 之间的关系，并给出证明; (3) 当 \(k > 2\) 时，求证: 在区间 \([-1, 5]\) 上，\(y = k(x + 3)\) 的图象位于函数 \(f(x)\) 图象的上方.
\end{problem}

\begin{problem}
已知数列 \(a_1, a_2, \cdots, a_{30}\)，其中 \(a_1, a_2, \cdots, a_{10}\) 是首项为 1，公差为 1 的等差数列: \(a_{10}, a_{11}, \cdots, a_{20}\) 是公差为 \(d\) 的等差数列: \(a_{20}, a_{21}, \cdots, a_{30}\) 是公差为 \(d^2\) 的等差数列 (\(d \neq 0\)). (1) 若 \(a_{20} = 40\)，求 \(d\); (2) 试写出 \(a_{30}\) 关于 \(d\) 的关系式，并求 \(a_{30}\) 的取值范围; (3) 续写已知数列，使得 \(a_{30}, a_{31}, \cdots, a_{40}\) 是公差为 \(a^{13}\) 的等差数列，\(\cdots, \cdots\)，依次类推，把已知数列推广为无穷数列. 提出同 (2) 类似的问题，(2) 应当作为特例，并进行研究，你能得到什么样的结论?
\end{problem}

\begin{problem}
(第1至12题)
\end{problem}

\begin{problem}
\(\frac{3}{4}\) 2.2. 3. \(\frac{1}{3} (x - 5), x \in [5,8]\).
\end{problem}

\begin{problem}
\( \left(-1,\frac{1}{2}\right). \) 5. \( (0,\sqrt{10}). \)
\end{problem}

\begin{problem}
-x - x\( ^{4} \). 7. 48.
\end{problem}

\begin{problem}
\(\frac{16}{3}\). 9. \(\frac{7}{25}\). 10. 2. 11. 4.
\end{problem}

\begin{problem}
\(\frac{a_1 + a_2 + \cdots + a_m}{m} \leqslant \frac{a_1 + a_2 + \cdots + a_n}{n}\)
(1 \( \leqslant m < n \) ) 和 \( \frac{a_{m+1} + a_{m+2} + \cdots + a_{n}}{n - m} \) \( \geqslant \frac{a_{1} + a_{2} + \cdots + a_{n}}{n} \) ( \( 1 \leqslant m < n \) ).
\end{problem}

\begin{problem}
(第13至16题) 题号 13 14 15 16 代号 B C A B
\end{problem}

\begin{problem}
(第17至22题)
\end{problem}

\begin{problem}
解法一 连接 \(A_{1}D, \text{因为} A_{1}D // B_{1}C\)， 因此 \( \angle BA_{1}D \) 为异面直线 \( A_{1}B \) 与 \( B_{1}C \) 所成的角. 连接 \(BD, D\) 在 \(\triangle A_{1}DB\) 中，\(A_{1}B = A_{1}D = 5\)，\(BD = 4\sqrt{2}\)，则 \[ \begin{array}{l} \cos \angle B A _{1} D = \frac {A _{1} B ^{2} + A _{1} D ^{2} - B D ^{2}}{2 \cdot A _{1} B \cdot A _{1} D} = \\ \frac {2 5 + 2 5 - 3 2}{2 \cdot \frac {5 \cdot 5}{2}} = \frac {9}{2 5}. \end{array} \] 因此 异面直线 \( \overrightarrow{A_{1}B} \) 与 \( B_{1}C \) 所成角的大小为 \( \arccos \frac{25}{25} \) .
解法二以 D 为坐标原点，分别以 DA、DC、 \( DD_{1} \) 所在直线为 x 轴，y 轴，z 轴，建立空间直角坐标系. 则 \[ A _{1} (4, 0, 3), B (4, 4, 0), B _{1} (4, 4, 3), C (0, 4, 0), \] 得 \( \overrightarrow{A_{1}B}=(0,4,-3),\overrightarrow{B_{1}C}=(-4,0,-3). \) 设 \(\overrightarrow{A_1B}\) 与 \(\overrightarrow{B_1C}\) 的夹角为 \(\theta\) 则 \(\cos \theta = \frac{\overrightarrow{A_1B}\cdot\overrightarrow{B_1C}}{|\overrightarrow{A_1B}|\cdot|\overrightarrow{B_1C}|} = \frac{9}{25}\)， 因此 \( \overrightarrow{A_{1}B} \) 与 \( \overrightarrow{B_{1}C} \) 的夹角大小为 \( \arccos\frac{9}{25} \) ，即异面直线 \(A_{1}B\) 与 \(B_{1}C\) 所成角的大小为 \(\arccos \frac{9}{25}\).
\end{problem}

\begin{problem}
解法一 \( \text{因为} w(1+2\i)=4+3\i, \) \[ \text{所以} w = \frac {4 + 3 \i}{1 + \frac {5}{2} \i} = 2 - \i, \] \[ \text{所以} z = \frac {2}{2 - 1} + | - \i | = 3 + \i. \] 若实系数一元二次方程有虚根 \(z = 3 + \i\) ，则必有其iden虚根 \(\bar{z} = 3 - \i\) \[ \text{因为} z + \bar {z} = 6, z \cdot \bar {z} = 1 0, \] 因此 所求的一个一元二次方程可以是 \( x^{2}-6x+10=0 \) . [解法二] 设 \(w = a + \mathrm{bi}, (a, b \in \R)\) \[ a + b \i - 4 = 3 \i - 2 a \i + 2 b, \] 得 \( \left\{\begin{aligned}a-4&=2b,\\ b&=3-2a,\end{aligned}\right. \) \( \text{所以}\left\{\begin{aligned}a&=2,\\ b&=-1,\end{aligned}\right. \) \[ \text{所以} w = 2 - \i, \] 以下解法同解法一].
\end{problem}

\begin{problem}
[解](1) \(\text{因为} \sin x = \frac{4}{5}, x \in \left[\frac{\pi}{2}, \pi\right]\)，\[ \text{所以} \cos x = - \frac {3}{5}, \] \[ f (x) = 2 \left(\frac {\sqrt {3}}{2} \sin x + \frac {1}{2} \cos x\right) - 2 \cos x \] \[ = \sqrt {3} \sin x - \cos x \] \[ = \frac {3}{5} \sqrt {3} + \frac {3}{5}. \] (2) \(f(x) = 2\sin \left(x - \frac{\pi}{6}\right)\)， \[ \text{因为} \frac {\pi}{2} \leqslant x \leqslant \pi , \] \[ \text{所以} \frac {\pi}{3} \leqslant x - \frac {\pi}{6} \leqslant \frac {5 \pi}{6} + \frac {1}{2} \leqslant \sin \left(x - \frac {\pi}{6}\right) \leqslant 1, \] 因此 函数 \( f(x) \) 的值域为 [1,2].
\end{problem}

\begin{problem}
[解](1) 设曲线方程为 \( y = ax^{2} + \frac{64}{7} \)，由题意可知， \(0 = a\cdot 64 + \frac{64}{7}.\text{所以} a = -\frac{1}{7}\). 因此 曲线方程为 \( y = -\frac{1}{7}x^{2} + \frac{64}{7} \) .
(2) 设变轨点为 \(C(x, y)\)，根据题意可知 \[ \left\{ \begin{array}{l l} \frac {x ^{2}}{1 0 0} + \frac {y ^{2}}{2 5} = 1, \& (1) \\ y = - \frac {1}{7} x ^{2} + \frac {6 4}{7}, \& (2) \end{array} \right. \] 得 \(4y^{2} - 7y - 36 = 0\) \(y = 4\) 或 \(y = -\frac{9}{4}\) （不合题意，舍去）. \[ \text{所以} y = 4. \] 得 \(x = 6\) 或 \(x = -6\) （不合题意，舍去）. 因此 C 点的坐标为 \( (6,4) \) ，\[ \vert A C \vert = 2 \sqrt {5}, \vert B C \vert = 4. \] 答: 当观测点 A、B 测得 AC、BC 距离分别为 \( 2\sqrt{5} \) 、4 时，应向航天器发出变轨指令.
\end{problem}

\begin{problem}
[解](1)
(2) 方程 \( f(x)=5 \) 的解分别是 \( 2-\sqrt{14},0,4 \) 和 \( 2+\sqrt{14} \) ，由于 \( f(x) \) 在 \( (-\infty,-1] \) 和 \( [2,5] \) 上单调递减，在 \( [-1,2] \) 和 \( [5,+\infty) \) 上单调递增，因此 \[ A = (- \infty , 2 - \sqrt {1 4} ] \cup [ 0, 4 ] \cup [ 2 + \sqrt {1 4}, + \infty). \] 由于 \(2 + \sqrt{14} < 6, 2 - \sqrt{14} > -2\)， \[ \text{所以} B \subset A. \] (3) [解法一] 当 \(x \in [-1, 5]\) 时， \[ f (x) = - x ^{2} + 4 x + 5. \] \[ g (x) = k (x + 3) - (- x ^{2} + 4 x + 5) \] \[ = x ^{2} + (k - 4) x + (3 k - 5) \] \[ = \left(x - \frac {4 - k}{2}\right) ^{2} - \frac {k ^{2} - 2 0 k + 3 6}{4}, \] \[ \text{因为} k > 2, \text{所以} \frac {4 - k}{2} < 1. \text {又} - 1 \leqslant x \leqslant 5, \] \circled{1} 当 \( -1 \leqslant \frac{4 - k}{2} < 1 \) ，即 \( 2 < k \leqslant 6 \) 时， 取 \( x = \frac{4 - k}{2} \) ， \( g(x)_{\min} = -\frac{k^{2} - 20k + 36}{4} = -\frac{1}{4}[(k - 10)^{2} - 64] \) . \[ \text{因为} 1 6 \leqslant (k - 1 0) ^{2} < 6 4, \] \[ \text{所以} (k - 1 0) ^{2} - 6 4 < 0, \] 则 \(g(x)_{\min} > 0\) \circled{2} 当 \( \frac{4-k}{2}<-1 \) ，即 k>6 时，取 x=-1， \( g(x)_{\min}=2k>0. \) 由\circled{1}、\circled{2}可知，当 \(k > 2\) 时，\(g(x) > 0, x \in [-1, 5]\). 因此，在区间 \([-1,5]\) 上，\(y = k(x + 3)\) 的图象位于函数 \(f(x)\) 图象的上方. [解法二] 当 \(x \in [-1, 5]\) 时， \[ f (x) = - x ^{2} + 4 x + 5. \] 由 \(\left\{ \begin{array}{l} y = k(x + 3), \\ y = -x^2 + 4x + 5, \end{array} \right.\) 得 \[ x ^{2} + (k - 4) x + (3 k - 5) = 0, \] 令 \(\Delta = (k - 4)^2 - 4(3k - 5) = 0\)，解得 \(k = 2\) 或 \(k = 18\) 在区间 \([-1,5]\) 上，当 \(k = 2\) 时，\(y = 2(x + 3)\) 的图象与函数 \(f(x)\) 的图象只交于一点 \((1,8)\)；当 \(k = 18\) 时，\(y = 18(x + 3)\) 的图象与函数 \(f(x)\) 的图象没有交点. 如图可知，由于直线 \( y = k(x + 3) \) 过点 \( (-3, 0) \) ，当 k > 2 时，直线 \( y = k(x + 3) \) 是由直线 \( y = 2(x + 3) \) 绕点 \( (-3, 0) \) 逆时针方向旋转得到. 因此，在区间 \( [-1, 5] \) 上，\( y = k(x + 3) \) 的图象位于函数 \( f(x) \) 图象的上方.
\end{problem}

\begin{problem}
[解] (1) \(a_{10} = 10\). \(a_{20} = 10 + 10d = 40\)，\(\text{所以} d = 3\). \[ \begin{array}{l} (2) a _{3 0} = a _{2 0} + 1 0 d ^{2} \\ = 1 0 (1 + d + d ^{2}) (d \neq 0), \\ a _{3 0} = 1 0 \left[ \left(d + \frac {1}{2}\right) ^{2} + \frac {3}{4} \right], \\ \text {当} d \in (- \infty , 0) \cup (0, + \infty) \text {时}, \\ a _{3 0} \in \left[ \frac {1 5}{2}, + \infty\right). \\ \end{array} \] (3) 所给数列可推广为无穷数列 \(\{a_{n}\}\)，其中 \(a_{1}, a_{2}, \cdots, a_{10}\) 是首项为 1，公差为 1 的等差数列，当 \(n \geqslant 1\) 时，数列 \(a_{10n}, a_{10n+1}, \cdots, a_{10(n+1)}\) 是公差为 \(d^{n}\) 的等差数列. 研究的问题可以是：试写出 \(a_{10(n + 1)}\) 关于 \(d\) 的关系式，并求 \(a_{10(n + 1)}\) 的取值范围. 研究的结论可以是：由 \(a_{40} = a_{30} + 10d^3 = 10(1 + d + d^2 +d^3)\) ，依次类推可得 \(a_{10(n + 1)} =\) \(10(1 + d + \dots +d^{m}) = \left\{ \begin{array}{ll}1 - d^{n + 1}\\ 10(n + 1), & d\neq 1,\\ 0(n + 1), & d = 1. \end{array} \right.\) 当 \(d > 0\) 时，\(a_{10(n + 1)}\) 的取值范围为 \((10, +\infty)\) 等.
\end{problem}

